\chapter{Results}
\chaptermark{} %this is the chapter heading that will show on subsequent pages

\section{Prediction Skill for Synthetic Data}

I first present the results pertaining to the accurarcy of forecasts for synthetic data.

\subsection{Sampling frequency, distribution}

Increasing sampling frequency led to improved forecast accurary, up to a finite cutoff where the error covariance inflation requirement becomes unstable.

Observational network density also improves forecast accuracy, as would be expected, up to limit of the computational memory avaiable for computing large non-sparse covariance matrices.
For example the full system, with 216,000 variables, produces an error covariance matrix of 347GB in double precision.
With the expected continuation of Moore's Law, however, this simulation may be possible in a few as 5 years (assuming a doubling of memory every 18 months, starting now with my 8GB lapop) on a laptop computer.

\subsection{Random error vs PMD perturbation scheme}

The computational implementation of the NS equations is dispersive, meaning that small perturbations are quickly resolved.

In other words, we beat the hell of out random error.

\subsection{Maintaining hydrostatic balance}

The importance of maintaining balance in the model variables, originally dooming the efforts of Richardson, was taken into consideration.

\section{Real-time Flow Reversal Prediction Skill}

Now the real business: running CFD in real time and assimilating observations from temperature probes.


